% Chapter 4

\chapter{System Architecture and Design} % Main chapter title

\label{Chapter4} % For referencing the chapter elsewhere, use \ref{Chapter4}

%----------------------------------------------------------------------------------------

In this chapter, we walk through the structural design of the IIITA Timetabling System. We cover the overall application topology, the relational database schema, how we ingest data from Excel files, the solver module architecture, and the routing logic that ties everything together.

%-----------------------------------
\section{Application Topology}

Our system follows a classic three-tier architecture adapted for the modern web using a Backend-as-a-Service (BaaS) model.

\begin{itemize}
    \item \textbf{Client Tier (React SPA):} The frontend is a Single Page Application (SPA) responsible for all user interactions, data formatting, client-side clash detection, and PDF generation. It communicates asynchronously with the backend via secure REST API calls over HTTPS.
    \item \textbf{Backend Tier (Supabase):} Acting as the serverless backend, Supabase provides an auto-generated API gateway, handles user authentication, and interfaces directly with the underlying database.
    \item \textbf{Data Tier (PostgreSQL):} The relational database stores all normalized institutional data, including users, entities (rooms, subjects, professors), and the core timetable slots.
\end{itemize}

%-----------------------------------
\section{Database Schema Design}

A well-structured relational schema is paramount for a scheduling system to maintain data integrity and support complex queries. The database is heavily normalized to prevent data duplication and anomalies. The core schema comprises the following primary entities and relationships:

\subsection{Core Entities}

\begin{figure}[htbp]
\centering
\begin{tikzpicture}[
    node distance=1.5cm and 2cm,
    entity/.style={rectangle, draw, fill=blue!10, rounded corners, text width=3cm, align=center, minimum height=1cm, font=\sffamily\small},
    relation/.style={diamond, draw, fill=orange!10, aspect=2, text width=2cm, align=center, font=\sffamily\scriptsize},
    line/.style={draw, thick, -{Latex[length=2mm]}},
    every node/.style={font=\sffamily}
]

% Nodes
\node[entity] (slots) {\textbf{timetable\_slots}\\ \scriptsize (id, start, end, day)};
\node[entity, above=of slots] (timetables) {\textbf{timetables}\\ \scriptsize (id, semester, status)};
\node[entity, left=of slots] (subjects) {\textbf{subjects}\\ \scriptsize (code, name, credits)};
\node[entity, right=of slots] (professors) {\textbf{professors}\\ \scriptsize (email, name, dept)};
\node[entity, below left=of slots] (rooms) {\textbf{rooms}\\ \scriptsize (name, capacity)};
\node[entity, below right=of slots] (groups) {\textbf{student\_groups}\\ \scriptsize (name, semester)};

% Relationships
\draw[line] (slots) -- (timetables) node[midway, left] {belongs to};
\draw[line] (slots) -- (subjects) node[midway, above] {teaches};
\draw[line] (slots) -- (professors) node[midway, above] {assigned to};
\draw[line] (slots) -- (rooms) node[midway, right, xshift=2mm] {hosted in};
\draw[line] (slots) -- (groups) node[midway, left, xshift=-2mm] {attended by};

\end{tikzpicture}
\caption{Simplified Entity-Relationship Diagram of the Core Schema}
\label{fig:er_diagram}
\end{figure}

\begin{itemize}
    \item \textbf{`timetables`:} Represents a specific iteration of a schedule (e.g., "B.Tech Sem 3 - Monsoon 2024"). It acts as the parent container for all scheduled slots. Attributes include the academic year, semester, publication status, and institutional break times (lunch start/end).
    \item \textbf{`subjects`:} Contains metadata for all offered courses. Attributes include the unique subject code, name, credit distribution (Lectures, Tutorials, Practicals), subject type (Core, Elective, Minor), and elective grouping data.
    \item \textbf{`professors`:} Stores faculty information, uniquely identified by email, alongside their departmental affiliation.
    \item \textbf{`rooms`:} Catalogs the physical infrastructure of the campus. Attributes include the unique room name (e.g., "CC-3-1001"), seating capacity, and room type (e.g., Lecture Hall, Computer Lab).
    \item \textbf{\texttt{student\_groups}:} Represents a specific cohort of students, defined by a unique combination of group name (e.g., "Sec A") and semester.
\end{itemize}

\subsection{The Central Hub: \texttt{timetable\_slots}}
The \texttt{timetable\_slots} table is the central junction of the schema. It represents a specific class occurrence in time and space. It acts as a massive join table, linking a specific \texttt{timetable\_id} to a \texttt{subject\_id}, \texttt{professor\_id}, \texttt{room\_id}, and \texttt{student\_group\_id}. 

Crucially, it also defines the temporal coordinates of the event, as detailed in Table~\ref{tab:schema_slots}.

\begin{table}[htbp]
\centering
\caption{Schema Specification for \texttt{timetable\_slots}}
\label{tab:schema_slots}
\renewcommand{\arraystretch}{1.4}
\begin{tabularx}{\textwidth}{@{} l l X @{}}
\toprule
\textbf{Column Name} & \textbf{Data Type} & \textbf{Description \& Constraints} \\
\midrule
\texttt{id} & \texttt{UUID} & Primary Key, auto-generated. \\
\texttt{timetable\_id} & \texttt{UUID} & Foreign Key referencing \texttt{timetables}. \\
\texttt{subject\_id} & \texttt{UUID} & Foreign Key referencing \texttt{subjects}. \\
\texttt{professor\_id} & \texttt{UUID} & Foreign Key referencing \texttt{professors}. \\
\texttt{room\_id} & \texttt{UUID} & Foreign Key referencing \texttt{rooms}. \\
\texttt{day\_of\_week} & \texttt{Integer} & $1 = \text{Monday}, \dots, 5 = \text{Friday}$. \\
\texttt{start\_time} & \texttt{VarChar} & 24-hour format (e.g., "09:00"). \\
\texttt{end\_time} & \texttt{VarChar} & 24-hour format (e.g., "11:00"). \\
\texttt{slot\_type} & \texttt{VarChar} & Enum: 'Lecture', 'Tutorial', 'Practical'. \\
\bottomrule
\end{tabularx}
\end{table}

\subsection{Generator and Solver Support Entities}
To support automated timetable generation, the schema includes mapping tables that define the constraints and requirements for specific semesters. These tables are actively used by the Generator Wizard and the solver engine:
\begin{itemize}
    \item \textbf{\texttt{semester\_clusters}:} Defines active academic clusters (e.g., Batch 2021, Semester 5, IT Department). The generator UI uses these clusters to scope the scheduling problem.
    \item \textbf{\texttt{cluster\_requirements}:} A many-to-many join table linking a \texttt{semester\_cluster} to the \texttt{subjects} that must be scheduled for that cohort. Includes \texttt{estimated\_enrollment} for room utilization optimization.
    \item \textbf{\texttt{professor\_expertise}:} A mapping table linking \texttt{professors} to \texttt{subjects} they are qualified to teach, along with a preference level. This is used by the auto-fill feature in the Generator Wizard to intelligently suggest professor assignments.
\end{itemize}

\subsection{Solver Data Model: Class Diagram}

Figure~\ref{fig:class_diagram} presents a UML class diagram of the solver's core data model, showing the relationships between the key TypeScript interfaces that drive the scheduling engine.

\begin{figure}[htbp]
\centering
\resizebox{\textwidth}{!}{%
\begin{tikzpicture}[
    class/.style={
        rectangle split, rectangle split parts=3,
        draw, thick, fill=blue!5,
        text width=5cm, align=center,
        font=\sffamily\scriptsize,
        rectangle split part fill={blue!15, white, white}
    },
    comp/.style={draw, thick, ->, >=Diamond[open], line width=0.8pt},
    assoc/.style={draw, thick, ->, >=stealth, dashed},
    note/.style={rectangle, draw, dashed, fill=yellow!10, text width=3cm, align=center, font=\sffamily\tiny, rounded corners}
]

% ClassSession
\node[class] (cs) {
    \textbf{ClassSession}
    \nodepart{second}
    \begin{tabular}{l}
    + id: number \\
    + subjectId: string \\
    + groupId: string \\
    + professorId: string \\
    + duration: number \\
    + slotType: enum \\
    + homeRoomIndex: number \\
    + isElective: boolean \\
    + basketName: string? \\
    + isLocked: boolean \\
    + lecturePairIndex: number \\
    \end{tabular}
    \nodepart{third}
    \textit{(immutable metadata)}
};

% Gene
\node[class, right=3cm of cs] (gene) {
    \textbf{Gene}
    \nodepart{second}
    \begin{tabular}{l}
    + day: 1--5 \\
    + startBucket: 1--8 \\
    + roomIndex: number \\
    \end{tabular}
    \nodepart{third}
    \textit{(scheduling decision)}
};

% SolverInput
\node[class, below=2.5cm of cs] (si) {
    \textbf{SolverInput}
    \nodepart{second}
    \begin{tabular}{l}
    + sessions: ClassSession[] \\
    + rooms: Room[] \\
    + numDays: 5 \\
    + numBuckets: 8 \\
    + config: SolverConfig \\
    \end{tabular}
    \nodepart{third}
    \textit{(solver entry point)}
};

% FitnessResult
\node[class, below=2.5cm of gene] (fr) {
    \textbf{FitnessResult}
    \nodepart{second}
    \begin{tabular}{l}
    + total: number \\
    + hardViolations: number \\
    + gapPenalty: number \\
    + roomUtilPenalty: number \\
    + profSameHalfPenalty: number \\
    + violationBreakdown: {...} \\
    \end{tabular}
    \nodepart{third}
    \textit{(evaluation output)}
};

% Solution type alias note
\node[note, above=0.8cm of gene] (sol) {
    \textbf{Solution} = Gene[] \\
    (1:1 with sessions)
};

% Relationships
\draw[comp] (si) -- node[left, font=\sffamily\tiny] {contains *} (cs);
\draw[comp] (sol) -- node[right, font=\sffamily\tiny] {contains *} (gene);
\draw[assoc] (cs) -- node[above, font=\sffamily\tiny] {maps 1:1} (gene);
\draw[assoc] (si.east) -- node[below, font=\sffamily\tiny] {evaluate() $\to$} (fr.west);

\end{tikzpicture}
}%
\caption{UML Class Diagram of the Solver Data Model}
\label{fig:class_diagram}
\end{figure}

%-----------------------------------
\section{The Data Ingestion Pipeline}

A major architectural feature of the system is its ability to ingest unstructured data. The transition from an administrator's raw Excel file to normalized database records follows a strict pipeline managed by the \texttt{TimetableImporter} module:

\begin{figure}[htbp]
\centering
\begin{tikzpicture}[
    node distance=1.2cm,
    block/.style={rectangle, draw, fill=gray!10, text width=10cm, align=center, rounded corners, minimum height=0.8cm, font=\sffamily\small},
    arrow/.style={draw, thick, ->, >=stealth}
]

\node[block] (step1) {\textbf{1. Binary Parsing}\\Reads \texttt{.xlsx} via SheetJS and converts to 2D JSON Array};
\node[block, below=of step1] (step2) {\textbf{2. Header Synchronization}\\Scans top rows to extract 1-hour and multi-hour time bounds};
\node[block, below=of step2] (step3) {\textbf{3. Metadata Extraction}\\Uses Regex to map subject codes, credits, and complex faculty assignments};
\node[block, below=of step3] (step4) {\textbf{4. Grid Traversal}\\Iterates over cells to infer class type (L/T/P) and extract room details};
\node[block, below=of step4] (step5) {\textbf{5. Review and Mutation}\\Displays editable React data grid; fires asynchronous \texttt{UPSERT} commands to Supabase};

\draw[arrow] (step1) -- (step2);
\draw[arrow] (step2) -- (step3);
\draw[arrow] (step3) -- (step4);
\draw[arrow] (step4) -- (step5);

\end{tikzpicture}
\caption{Data Ingestion Pipeline for Excel Timetables}
\label{fig:parsing_pipeline}
\end{figure}

\begin{enumerate}
    \item \textbf{Binary Parsing:} The user uploads an `.xlsx` file. The `xlsx` library reads the binary stream and converts the active spreadsheet into a two-dimensional JSON array.
    \item \textbf{Header Synchronization:} The heuristic engine scans the top rows to identify the time boundary headers. It extracts standard 1-hour blocks (e.g., 9:00-10:00) and multi-hour blocks.
    \item \textbf{Metadata Extraction:} The engine locates the "Course Code / Faculties" section of the spreadsheet. It applies Regular Expressions (Regex) to map subject codes to full names and extracts the L-T-P-S credit structures. It also parses complex faculty assignment strings (e.g., identifying that "Prof. X" teaches "Sec A \& B" while "Prof. Y" teaches "Sec C").
    \item \textbf{Grid Traversal and Extraction:} The engine iterates over the main timetable grid. For each cell, it:
    \begin{itemize}
        \item Identifies the current day from the leftmost column.
        \item Identifies the time bounds from the column header.
        \item Extracts the subject code, infers the class type (L/T/P) based on cell coloring or text markers, and extracts room and section assignments.
    \end{itemize}
    \item \textbf{Review and Mutation:} The extracted raw data is presented to the administrator in an editable React data grid. Upon confirmation, the application fires a sequence of asynchronous \texttt{UPSERT} commands to Supabase, ensuring that all necessary entities (professors, rooms, subjects) exist before performing a bulk \texttt{INSERT} into the \texttt{timetable\_slots} table.
\end{enumerate}

%-----------------------------------
\section{View Routing and State Management}

The application does not utilize a traditional URL-based router (like React Router). Instead, it employs state-based view switching managed by the root `App.tsx` component. 

A central state variable, `currentView`, dictates which heavy component is mounted into the main viewing area. The available views are:
\begin{itemize}
    \item \texttt{import}: Mounts the \texttt{TimetableImporter}.
    \item \texttt{generator}: Mounts the multi-step \texttt{GeneratorView} wizard for automated timetable creation.
    \item \texttt{editor}: Mounts the interactive \texttt{EditorView} for drag-and-drop timetable refinement with LNS auto-fix.
    \item \texttt{student}: Mounts the \texttt{TimetableViewer} (All Students administrative view).
    \item \texttt{prof}: Mounts the \texttt{ProfessorTimetableViewer}.
    \item \texttt{room}: Mounts the \texttt{RoomTimetableViewer}.
    \item \texttt{master-room}: Mounts the campus-wide \texttt{MasterRoomViewer}.
    \item \texttt{free-room}: Mounts the \texttt{FreeRoomViewer}.
    \item \texttt{report-card}: Mounts the \texttt{ReportCardView}.
\end{itemize}

This architecture ensures that only the data required for the active view is fetched from the database, minimizing unnecessary network traffic. Custom React hooks, notably \texttt{useGeneratorData} and \texttt{useEditorData}, manage the fetching and caching of complex relational data required by these components, abstracting the Supabase query syntax away from the presentation layer.
